\section{河西、城寨与文书所见的资源}

1009年李德明经营夏州、灵州，1036年党项势力扩向甘、凉，1044年宋夏和议，1082年永乐城失守：这些节点共同说明西夏的资源并不只来自一座王都。河西道路、水源、城寨、贸易与佛寺把不同区域相连；战争使这些联系更可见，却不能由战报直接量化税收、人口或粮运。

黑水城文书使行政、法律、借贷和寺院经济有机会进入叙事，但其出土集中性恰要求谨慎。文书所见的一个机构、一个户类或一次交易不能代表宁夏与河西全境。将宋方的战报、金辽外交记录、西夏文字材料和城址分期对读，才能把“边疆国家”的财政理解为地方道路、书写技术和资源调度的组合，而非朝贡关系的附属物。
